\section{Related Work}


% \subsection{Test-time Scaling}
% Recent studies on test-time scaling has progressed from simple prompt-engineering tricks to sophisticated, multi-stage systems that modulate inference compute on the fly. Early studies such as Chain-of-Thought prompting \citep{wei2023chainofthoughtpromptingelicitsreasoning} showed that instructing a model to verbalize intermediate reasoning can unlock latent capabilities, while Self-Consistency~\citep{wang2022self} and Best-of-N sampling demonstrated that sampling multiple independent traces and voting improves reliability. Sequential strategies like STaR \citep{zelikman2022star} and Self-Refine \citep{madaan2023self} add iterative revision, and hybrid search methods—including Tree-of-Thoughts \citep{yao2024tree} and MCTS-based planners \citep{zhang2024o1}—combine breadth and depth to explore diverse reasoning paths under a compute budget. 

% Parallel trends~\citep{xie2025teaching} in internal scaling fine-tune models to decide autonomously how much “thinking” to allocate, as in OpenAI o1~\citep{jaech2024openai} and DeepSeek-R1~\citep{deepseekai2025deepseekr1incentivizingreasoningcapability}, often by pairing supervised traces with reinforcement-learning feedback. Collectively, these lines of work establish a taxonomy that frames test-time scaling as a unifying lens for improving LLM reasoning without changing model parameters at training time.

% Reinforcement Learning (RL) has demonstrated its significance for enhancing the capabilities of LLMs and LMMs. Early research primarily focused on Reinforcement Learning from Human Feedback (RLHF)~\citep{ouyang2022training} to align with human preferences. Recent work ~\citep{deepseekai2025deepseekr1incentivizingreasoningcapability} introduces the open-source DeepSeek-R1 model and investigates how RL with verifiable rewards can promote advanced chain-of-thought reasoning and reflective behaviors. Several recent studies have explored R1 a



\subsection{Multimodal Chain-of-Thought}
Chain-of-Thought (CoT) prompting was first shown to elicit step-by-step reasoning in LLMs by supplying a few worked examples that include intermediate rationales~\citep{feng2023towards,zhang2024chain, wei2023chainofthoughtpromptingelicitsreasoning}.
Recent extensions of CoT to multimodal settings embed visual evidence directly into the reasoning trajectory. ICoT ~\citep{zhang2024multimodalchainofthoughtreasoninglanguage} interleaves attention-selected image crops with text tokens, yielding significant VQA gains, while Visual CoT ~\citep{shao2024visualcotadvancingmultimodal} supplies 438 k bounding-box-grounded rationales to train VLMs that emit explicit visual tokens and improve spatial grounding.
Recent works~\citep{hu2024visual,zhou2024image,yang20253d,gao2024cantor, wu2025reinforcing, chern2025thinking,fang2025intention, cheng2025visual, su2025openthinkimg} further leverage external tools to supply visual cues that enrich multimodal CoT reasoning..

Recent works~\citep{chen2025blip3, wang2025autoregressive} like Chameleon~\citep{chameleonteam2025chameleonmixedmodalearlyfusionfoundation,chern2024anole} trains a unified token-based model that can emit arbitrary sequences of text and image tokens, but at the cost of large-scale pixel-level supervision and heavier decoding. MVoT~\citep{li2025imagine} further trains a unified model to directly produce image and text interleaving trajectories, but absent of reasoning thoughts.
In contrast, our \Model framework differs by emitting compact latent visual tokens rather than real image patches or pixels, avoiding heavy image generation while still allowing fully interleaved visual–text reasoning.

\subsection{Latent Reasoning in LLMs}

Much recent work has highlighted the importance of intermediate hidden representations in Large Language Models (LLMs)~\citep{biran2024hopping,yang2024large}. To better guide the latent reasoning process, several approaches introduce specialized tokens into the input sequence. ~\cite{wang2023guiding} incorporate discrete \texttt{<plan>} tokens to control reasoning stages, while ~\cite{goyal2023think} propose inserting a \texttt{<pause>} token during pretraining to stabilize multi-step reasoning.

Another line of work seeks to internalize reasoning behavior by distilling chain-of-thought rationales into latent representations. ~\cite{deng2023implicit} trains models to mimic CoT-style reasoning implicitly through hidden states, and ~\citep{deng2024explicit} further improves inference efficiency by removing explicit intermediate steps altogether. ~\cite{yu2024distilling} proposes to distill latent reasoning capabilities into a model by supervising it with data generated for complex reasoning. More recently, ~\cite{hao2024training} go further by replacing CoT tokens with continuous latent embeddings, enabling unconstrained reasoning in the latent space to explore on complex tasks including math and logical reasoning.
While prior work primarily focuses on enhancing efficiency or structural planning within the LLM's latent space, our approach takes a different perspective: we treat latent tokens as a \emph{bridge for exploring visual information} into the model.

% Rather than confining reasoning within the language model’s internal space, we expand the latent trajectory to incorporate cross-modal signals, enabling the model to reason more effectively on visually grounded tasks. This shifts the objective from internal reasoning compression to multimodal integration and alignment.
